\section{Introduction}

\subsection{Large Language Models (LLMs)}

\subsection{Software 3.0}
The term "Software 3.0," introduced by Andrej Karpathy, describes how AI is reshaping the way we build software.
\begin{itemize}
    \item \textbf{Software 1.0}: explicit rules written by programmers (if/else logic, loops).
    \item \textbf{Software 2.0}: behavior learned from data (neural networks trained on examples).
    \item \textbf{Software 3.0}: behavior specified in natural language (LLMs following prompts).
\end{itemize}

\subsection{Roles of AI}
Over the past decade, AI has moved from a niche technology to a core part of software design. In practice, AI systems tend to play three main roles:
\begin{itemize}
    \item \textbf{AI as a Coding Agent}: AI systems can write, refactor, and review code, accelerating software development. 
    Their output is largely deterministic because it is traditional code that can be inspected, tested, and debugged.
    \item \textbf{AI as a Component}: In this case, AI is embedded within a larger system, and the surrounding architecture is designed around its capabilities and limitations. 
    This is not entirely new: many traditional machine learning models already fit this pattern.
    \item \textbf{AI as a User}: AI systems can also act as users of other systems, because they can consume inputs and generate outputs through your interfaces.
    Today, system design must account for AI clients that can call APIs, read documentation, and operate on application data.
    This is introducing a new question: Who are you building for? Human users, AI users, or both?
\end{itemize}
To design responsibly, we should distinguish when AI is unnecessary, when it is simply convenient, and when it is a true system component.
Those choices are important since they respond to the question "What are you building for?" and this mainly depends on how reniewable and fixed we want our outputs.
\begin{itemize}
    \item \textbf{No AI}: Use traditional software when rules are clear, logic is straightforward, and outputs are deterministic and easy to test.
    \item \textbf{AI as a Convenience}: The problem and rules are well understood, but prompting is faster or simpler than writing code.
    Outputs are still mostly predictable and testable, but reliability checks remain important.
    \item \textbf{AI as a System Component}: Use AI when behavior is too complex or context-dependent to encode with explicit rules.
    In these cases, conventional testing is harder because expected outputs are open-ended, and the value comes from handling ambiguity (for example, creative generation and open-ended Q\&A).
\end{itemize}
This course focuses primarily on AI as a system component, while still covering the other two cases.
In this context, we want to analyze how reviewable and fixed outputs are, and for which kind of users we are building for.

\subsection{Challenges}
Creating AI systems is more complex than traditional software because of the following challenges:
\begin{itemize}
    \item \textbf{Getting a single call right}: LLMs are not deterministic, and their outputs can vary widely based on the prompt, which needs to be accurately designed to get the desired output.
    \item \textbf{Managing a conversation}: When building a system that interacts with an LLM, you need to manage the context across the conversation, which can make the costs grow rapidly.
    \item \textbf{Long-term memory}: LLMs have limited context windows, so you need to design strategies for storing and retrieving information (like preferences or conversation history) over time.
    \item \textbf{Enterprise/Community knowledge}: LLMs could use private data to train, so we need to design systems that can leverage this knowledge while respecting privacy and security concerns, also of our enterprise.
\end{itemize}